% !TEX root = ../algebra_02.tex

\section{Гомоморфизм, образ и ядро}

\begin{Definition}{Гомоморфизм групп}{def:hom}
	Пусть $G$ и $G'$ --- группы. Отображение $\varphi\colon G \to G'$
	называется \emph{гомоморфизмом групп}, если для всех $g_1,g_2\in G$ выполнено
	\[
	\varphi(g_1g_2)=\varphi(g_1)\varphi(g_2).
	\]
\end{Definition}

\begin{Example}{}{ex:hom-examples}
	\begin{enumerate}
		\item Определитель:
		\[
		\det\colon \GL_n(K) \to K^{\times}, \qquad \det(AB)=\det(A)\det(B).
		\]
		\item Знак перестановки:
		\[
		\sgn\colon S_n \to \{\pm 1\}, \qquad \sigma \mapsto \sgn(\sigma).
		\]
		\item Тождественный и нулевой гомоморфизмы:
		\[
		\id_G\colon G\to G, \quad g\mapsto g;
		\qquad
		0\colon G\to G', \quad g\mapsto e_{G'}.
		\]
		\item Логарифм как гомоморфизм между мультипликативной и аддитивной группами:
		\[
		\ln\colon (\RR_{>0},\cdot) \to (\RR,+), \qquad a\mapsto \ln a.
		\]
	\end{enumerate}
\end{Example}

\begin{Lemma}{Сохранение единицы и обратного элемента}{lem:unit-inverse}
	Пусть $\varphi\colon G\to G'$ --- гомоморфизм. Тогда
	\begin{enumerate}
		\item $\varphi(e_G)=e_{G'}$;
		\item для любого $g\in G$ выполнено
		\[
		\varphi(g^{-1})=\varphi(g)^{-1}.
		\]
	\end{enumerate}
\end{Lemma}

\begin{proof}
	Так как $e_Ge_G=e_G$, имеем
	\[
	\varphi(e_G)\varphi(e_G)=\varphi(e_G).
	\]
	Умножая это равенство слева на $\varphi(e_G)^{-1}$, получаем $\varphi(e_G)=e_{G'}$.

	Далее,
	\[
	\varphi(g^{-1})\varphi(g)=\varphi(g^{-1}g)=\varphi(e_G)=e_{G'},
	\]
	и аналогично
	\[
	\varphi(g)\varphi(g^{-1})=e_{G'}.
	\]
	Следовательно, $\varphi(g^{-1})=\varphi(g)^{-1}$.
\end{proof}

\begin{Proposition}{Образы и прообразы подгрупп}{image-preimage}
	Пусть $\varphi\colon G\to G'$ --- гомоморфизм.
	\begin{enumerate}
		\item Если $H\le G$, то $\varphi(H)\le G'$.
		\item Если $H'\le G'$, то $\varphi^{-1}(H')\le G$.
		\item Если $H'\triangleleft G'$, то $\varphi^{-1}(H')\triangleleft G$.
	\end{enumerate}
\end{Proposition}

\begin{proof}
	\begin{enumerate}
		\item Пусть $g_1',g_2'\in \varphi(H)$. Тогда существуют $g_1,g_2\in H$ такие, что $g_1'=\varphi(g_1)$ и $g_2'=\varphi(g_2)$. Поскольку $H$ --- подгруппа, $g_1g_2\in H, \qquad g_1^{-1}\in H.$\\
		Следовательно, $g_1'g_2' = \varphi(g_1)\varphi(g_2)=\varphi(g_1g_2)\in \varphi(H),$ и по лемме $(g_1')^{-1}=\varphi(g_1)^{-1}=\varphi(g_1^{-1})\in \varphi(H).$\\
		Значит, $\varphi(H)\le G'$.
		\item Пусть $g_1,g_2\in \varphi^{-1}(H')$. Тогда $\varphi(g_1),\varphi(g_2)\in H'$.\\
		Так как $H'$ --- подгруппа, $\varphi(g_1g_2)=\varphi(g_1)\varphi(g_2)\in H',$ то есть $g_1g_2\in \varphi^{-1}(H')$.\\
		Кроме того, $\varphi(g_1^{-1})=\varphi(g_1)^{-1}\in H',$ поэтому $g_1^{-1}\in \varphi^{-1}(H')$. Значит, $\varphi^{-1}(H')\le G$.

		\item Пусть $g\in G$ и $h\in \varphi^{-1}(H')$. Тогда $\varphi(h)\in H'$. Поскольку $H'\triangleleft G'$, имеем
		\[
		\varphi(ghg^{-1})=\varphi(g)\varphi(h)\varphi(g)^{-1}\in H'.
		\]
		Следовательно, $ghg^{-1}\in \varphi^{-1}(H')$. Значит, $\varphi^{-1}(H')\triangleleft G$.
	\end{enumerate}
\end{proof}

\begin{Corollary}{Образ и ядро}{prop:image-kernel}
	Для любого гомоморфизма $\varphi\colon G\to G'$ имеем
	\[
	\Im \varphi := \varphi(G) \le G',
	\qquad
	\Ker \varphi := \varphi^{-1}(\{e_{G'}\}) = \{g\in G \mid \varphi(g)=e_{G'}\} \triangleleft G.
	\]
\end{Corollary}

\begin{proof}
	Первая часть следует из пункта 1) предложения~\ref{prop:image-preimage}, вторая --- из пункта 3) того же предложения, поскольку $\{e_{G'}\} \triangleleft G'$.
\end{proof}

\begin{Example}{}{ex:alternating}
	Поскольку знак перестановки $\sgn \colon S_n \to (\{ \pm 1 \}, \cdot)$ является гомоморфизмом, то его ядро
	\[
	A_n = \Ker(\sgn) \triangleleft S_n
	\]
	является нормальной подгруппой.
\end{Example}

\begin{Proposition}{Критерий инъективности}{prop:inj-kernel}
	Пусть $\varphi\colon G\to G'$ --- гомоморфизм. Тогда
	\[
	\varphi \text{ инъективен } \Longleftrightarrow \Ker\varphi = \{e_G\}.
	\]
\end{Proposition}

\begin{proof}
	Предположим, что $\varphi$ инъективен. Если $h\in \Ker\varphi$, то $\varphi(h)=e_{G'}=\varphi(e_G).$

	Из инъективности следует $h=e_G$, значит, $\Ker\varphi=\{e_G\}$.

	Обратно, пусть $\Ker\varphi=\{e_G\}$ и $\varphi(h_1)=\varphi(h_2).$

	Тогда $\varphi(h_1h_2^{-1})=\varphi(h_1)\varphi(h_2)^{-1}=e_{G'},$ то есть $h_1h_2^{-1}\in \Ker\varphi$. Следовательно, $h_1h_2^{-1}=e_G,$ откуда $h_1=h_2$. Значит, $\varphi$ инъективен.
\end{proof}

\begin{Definition}{Специальные виды гомоморфизмов}{def:special-homs}
	\begin{enumerate}
		\item Инъективный гомоморфизм называется \emph{мономорфизмом}.
		\item Сюръективный гомоморфизм называется \emph{эпиморфизмом}.
		\item Биективный гомоморфизм называется \emph{изоморфизмом}.
		\item Гомоморфизм $G\to G$ называется \emph{эндоморфизмом}.
		\item Изоморфизм $G\to G$ называется \emph{автоморфизмом}.
	\end{enumerate}
\end{Definition}

\begin{Proposition}{Обратное отображение к изоморфизму}{prop:inverse-iso}
	Если $\varphi\colon G\to G'$ --- изоморфизм, то обратное отображение
	\[
	\varphi^{-1}\colon G'\to G
	\]
	тоже является изоморфизмом.
\end{Proposition}

\begin{proof}
	Так как $\varphi$ биективно, отображение $\varphi^{-1}$ существует и биективно. Для любых $g_1',g_2'\in G'$ положим
	\[
	A:=\varphi^{-1}(g_1'g_2'), \qquad B:=\varphi^{-1}(g_1')\varphi^{-1}(g_2').
	\]
	Тогда
	\[
	\varphi(A)=g_1'g_2'=
	\varphi\bigl(\varphi^{-1}(g_1')\bigr)\,\varphi\bigl(\varphi^{-1}(g_2')\bigr)=\varphi(B).
	\]
	Из инъективности $\varphi$ следует $A=B$. Следовательно,
	\[
	\varphi^{-1}(g_1'g_2')=\varphi^{-1}(g_1')\varphi^{-1}(g_2').
	\]
	Значит, $\varphi^{-1}$ --- гомоморфизм, а потому изоморфизм.
\end{proof}
